\section{Modelo implementado}\label{sec:modelo}

\subsection{Organización general}

El modelo se ha implementado en Python 3 con la librería SimPy, como preveía la
Parte 1. Es una simulación terminante: cada réplica es un partido completo, que
empieza con el estadio vacío al abrir las puertas en $t=-90$~min y termina
cuando todos los espectadores han salido o están a salvo. No hay periodo de
calentamiento, porque el estadio vacío es el estado inicial real y la entrada
forma parte de las preguntas. El reloj de SimPy trabaja en segundos con $t=0$ en
el inicio del partido; los resultados se expresan en minutos respecto a ese
instante. El final ordinario del partido es $t=+105$~min (90 de juego y 15 de
descanso).

Todos los parámetros están en un fichero de configuración. Cada experimento
declara solo lo que cambia respecto al escenario base, lo que hace explícito el
principio de la Parte 1 de cambiar un solo factor. La carga de la configuración
la valida y rechaza claves desconocidas, de modo que un error tipográfico no
produce en silencio un escenario distinto del previsto. Ningún parámetro del
modelo está escrito en el código.

Cada espectador es un único proceso que encadena entrada, ubicación,
permanencia y salida; la alarma de emergencia puede interrumpirlo para pasar a
la evacuación. Las averías de los tornos son procesos auxiliares que se repiten
indefinidamente y no determinan el fin de la réplica. Al terminar, la réplica
devuelve una fila por espectador con todos sus instantes, y los indicadores se
calculan a partir de esas filas (sección~\ref{sec:indicadores}).

\subsection{Generación de espectadores y llegadas}

Al empezar la réplica se generan los $N=48.775$ espectadores con todos sus
atributos:

\begin{itemize}
  \item \textbf{Hora de llegada.} Se sortea la franja de 15~min con las
    probabilidades del perfil (base o anticipado, tabla de la Parte 1) y un
    instante uniforme dentro de ella.
  \item \textbf{Macrozona, puerta y bloque.} Macrozona con un 25~\% de
    probabilidad cada una y, dentro de ella, puerta (1--6) y bloque (1--10)
    uniformes. Los bloques 1--5 forman la grada inferior y los 6--10 la
    superior.
  \item \textbf{Grupo.} Estándar (88~\%), mayores (10~\%) o personas con
    movilidad reducida, PMR (2~\%), con un factor de velocidad de 1, 0,8 y 0,5.
  \item \textbf{Velocidad, distancia y servicio.} Velocidad Normal(1,45;
    0,23)~m/s truncada en [0,6;~2,2] y multiplicada por el factor del grupo;
    distancia de la puerta al asiento Uniforme(60,~180)~m, que se usa también
    para salir; tiempo de servicio en el torno Triangular(2,~4,~6)~s.
\end{itemize}

Los espectadores se ordenan por hora de llegada. Cada atributo sale de su
propia fuente de números aleatorios (apartado~\ref{sec:nac}).

\subsection{Entrada: colas, tornos y prioridad}

Cada puerta es un recurso con tantas plazas como tornos (3 en el escenario
base) y una cola con prioridades. El espectador espera hasta su hora de
llegada, elige puerta (apartado~\ref{sec:cierre}), espera un torno libre y lo
ocupa durante su tiempo de servicio y, si le toca, la relectura de la entrada.
La espera en cola es el inicio del control menos la llegada.

Con la disciplina FIFO todos piden con el mismo nivel. Con prioridad (X3),
mayores y PMR piden con un nivel más alto y adelantan a los demás en la cola de
su puerta. La prioridad no es expulsiva: nunca interrumpe un control ya
iniciado. Dentro de cada nivel se mantiene el orden de llegada. Las averías
tienen un nivel aún más alto.

\subsection{Cierre de puerta y redirección}\label{sec:cierre}

En el experimento X6 la puerta 1 de la macrozona A se cierra durante 10~min
(cierre de puerta, I4). Mientras está cerrada nadie entra por ella. La puerta
se elige una sola vez, al llegar, con el estado de ese instante:

\begin{itemize}
  \item si la puerta asignada está abierta, el espectador va a ella;
  \item si está cerrada y la redirección está activa, va a la puerta abierta de
    su macrozona con menos cola (se desempata por personas en servicio y, después,
    por número de puerta; es una regla determinista, sin azar);
  \item si está cerrada y no hay redirección, espera a que reabra y entra después
    por su puerta, en el mismo orden en que llegaron quienes esperaban.
\end{itemize}

Quien se redirige no recorre ninguna distancia adicional, porque el modelo no
representa el exterior del estadio. Quien ya estaba en la cola de la puerta
cuando se cierra sigue siendo atendido; es una hipótesis adoptada que se discute
en la sección~\ref{sec:desviaciones}.

\subsection{Incidencias en el torno}

\textbf{Fallo de lectura (I2).} Con probabilidad 3~\%, al terminar el servicio
la entrada se vuelve a leer durante Uniforme(5,~15)~s más, sin liberar el torno.
La relectura no interrumpe ni es interrumpida.

\textbf{Avería de torno (I3).} Cada torno alterna periodos de funcionamiento
Exponencial(media 240~min) y reparaciones Uniforme(3,~10)~min. El tiempo entre
fallos se mide en funcionamiento: el reloj se detiene durante la reparación, de
modo que 240~min es el tiempo medio de buen funcionamiento y un torno no puede
averiarse mientras está parado. Al vencer ese tiempo, la avería pide el torno con
la prioridad máxima y sin expulsión: el torno termina el servicio en curso y
después se para, como indicaba la Parte 1 («cuando el torno queda libre»). La
avería adelanta a toda la cola, incluidos los grupos priorizados.

\subsection{Ubicación en el asiento y conflictos}

Tras el torno, el espectador camina hasta el vomitorio de su bloque a su propia
velocidad. Cada bloque tiene un vomitorio de 8 plazas atendidas por orden de
llegada; el paso dura Triangular(5,~10,~20)~s y al salir del vomitorio se sienta.

\textbf{Conflicto por asiento (I1).} Al llegar a su asiento, el espectador lo
encuentra ocupado con probabilidad 0,5~\% en el escenario base y 1~\% en el
experimento de acomodadores (X5). Mientras dura el conflicto no suelta su plaza
del vomitorio: pide un acomodador de su macrozona (6 en el escenario base),
espera si están todos ocupados y el conflicto se resuelve en Triangular(1,~3,~10)~min
desde que llega el acomodador. Retener la plaza es lo que hace que un conflicto
retrase también al resto del bloque: con varios conflictos a la vez, el
vomitorio pierde plazas y quienes vienen detrás esperan más (pregunta P6).

\subsection{Salida ordinaria}

Al final del partido cada macrozona dispone de 6 salidas de 4~m. Con 66 personas
por metro y minuto, cada salida deja pasar 264 personas por minuto, una cada
0,227~s. Cada espectador vuelve a recorrer su distancia y, al llegar a la zona de
salidas, elige la de su macrozona con menos personas en ese momento. Se han
simulado dos políticas: salida simultánea en $t=+105$ y salida escalonada (X4),
en la que la grada superior sale 3~min después que la inferior.

\subsection{Evacuación de emergencia}

En X7 suena una alarma en $t=+40$~min. Se cierran los accesos, las salidas se
ensanchan un 30~\% (de 4 a 5,2~m, una persona cada 0,175~s) y cada espectador
actúa según dónde está:

\begin{itemize}
  \item quien aún no ha llegado ya no llega y no cuenta en la evacuación;
  \item quien está en la cola exterior o pasando el torno abandona y tampoco cuenta;
  \item quien ya ha pasado el torno, esté caminando, en el vomitorio, en un
    conflicto o sentado, evacúa: espera su tiempo de reacción,
    Triangular(30,~60,~120)~s, camina un 20~\% más rápido hasta la salida de su
    macrozona con menos cola y la atraviesa.
\end{itemize}

En la variante con bloqueo, una de las 6 salidas de la macrozona A queda
inutilizada y su público se reparte entre las otras 5.

\subsection{Números aleatorios comunes}\label{sec:nac}

La Parte 1 pedía usar las mismas semillas en todos los escenarios y un generador
independiente por fuente de azar. La implementación va un paso más allá, porque
las mismas semillas no bastan si los números se asignan en distinto orden:

\begin{itemize}
  \item \textbf{Un generador por fuente.} Cada una de las 13 fuentes de azar
    (llegadas, asignación, grupo, velocidad, distancia, servicio, relectura,
    vomitorio, conflicto, avería, elección de salida, reacción y desempate de
    puerta) tiene su propio generador, derivado de la semilla de la réplica y de
    la posición fija de la fuente. Consumir más o menos números de una fuente no
    altera las demás.
  \item \textbf{Todo sorteo por espectador se hace al empezar la réplica.} Si el
    servicio o la relectura se sortearan durante la simulación con un generador
    compartido, cada número caería en un espectador distinto según el orden en
    que el modelo lo atiende, y ese orden cambia entre escenarios (por ejemplo,
    con 2 o 4 tornos). Por eso el servicio, la relectura y su duración, el paso
    por el vomitorio, el conflicto y su resolución, el desempate entre salidas y
    el tiempo de reacción se sortean una vez por espectador y quedan asociados a
    él. Además, conflicto y relectura se deciden comparando un uniforme con la
    probabilidad: al pasar del 0,5~\% al 1~\% de conflictos se conservan los de
    la base y solo se añaden otros, y las tres variantes de X5 tienen
    exactamente los mismos conflictos.
  \item \textbf{Un generador por torno para las averías.} El calendario de averías
    de cada torno es el mismo con 2, 3 o 4 tornos por puerta y no depende de las
    demás macrozonas.
  \item \textbf{Semilla de cada réplica:} $20261105+r$ para la réplica $r$
    (numerada desde 0), igual en todos los escenarios.
\end{itemize}

Así, la diferencia entre dos escenarios con la misma semilla se debe solo a lo
que cambia entre ellos. El efecto se comprueba en el
apartado~\ref{sec:reproducibilidad}.
